% ===================================================================
%  Αρχείο: 37167_SOLUTION.tex (Fragment)
%  Αυτόματη μετατροπή μέσω doc2latex.py
% ===================================================================

ΘΕΜΑ 4

Δίνεται παραλληλόγραμμο ΑΒΓΔ με ΑΒ \textgreater{} ΑΔ και οι διχοτόμοι των γωνιών του ΑΡ, ΒΕ, ΓΣ και ΔΤ (όπου Ρ, Ε στην ΔΓ και Σ, Τ στην ΑΒ) τέμνονται στα σημεία Κ, Λ, M και N όπως φαίνεται στο παρακάτω σχήμα.

Να αποδείξετε ότι:

\begin{alist}
\item το τετράπλευρο ΔΕΒΤ είναι παραλληλόγραμμο. \monades{7}

\item το τετράπλευρο ΚΛΜΝ είναι ορθογώνιο. \monades{8}

\item ΛΝ // ΑΒ \monades{5}

\item ΛΝ = ΑΒ -- ΑΔ \monades{5}

\includesvg[width=5.52921in,height=2.63542in]{/home/spyros/Μαθηματικά/Trapeza_Thematwn/A_Lykeiou/Geometry/extracted/37167_SOLUTION/media/image2.svg}
\end{alist}
